\documentclass[aspectratio=169]{beamer}

% Theme and styling
\usetheme{metropolis}
\usepackage{booktabs}
\usepackage{graphicx}
\usepackage{subcaption}
\usepackage{amsmath}
\usepackage{xcolor}

% Title page information
\title{Thesis Update}
\subtitle{Evaluating Foundation Models for Financial Time-Series Forecasting}
\author{Dimitrios Kourntidis}
\institute{Leiden University}
\date{\today}

% Custom colors
\definecolor{LUBlue}{RGB}{0,48,87}
\definecolor{LUGreen}{RGB}{120,190,32}

\begin{document}

% Title slide
\begin{frame}
\titlepage
\end{frame}

% Outline
\begin{frame}{Outline}
\tableofcontents
\end{frame}

% ============================================================
\section{Research Questions}
% ============================================================

% \begin{frame}{Motivation}
% \begin{itemize}
% \item Financial time-series forecasting is critical for investment decisions
% \item Traditional statistical methods (ARIMA) vs. emerging foundation models
% \item \textbf{Key Question:} Can pretrained foundation models outperform traditional methods in stock price prediction?
% \end{itemize}

% \vspace{1em}

% \begin{block}{Research Gap}
% Limited comparative analysis of foundation models (Chronos, Sundial, TimesFM) against traditional baselines on financial data
% \end{block}
% \end{frame}

\begin{frame}{Research Questions}
\textbf{RQ1: Forecast Horizon Length Effects}
\begin{itemize}
\item How does forecast horizon length affect the relative performance of foundation models compared to traditional approaches?
\end{itemize}

\vspace{0.5em}

\textbf{RQ2: Temporal Frequency Effects}
\begin{itemize}
\item How does temporal frequency of input data affect the relative performance of foundation models versus traditional approaches?
\end{itemize}

\vspace{0.5em}

\textbf{RQ3: Market Disruption Resilience}
\begin{itemize}
\item How do foundation models compare to traditional approaches in maintaining accuracy during high-volatility market regimes?
\end{itemize}
\end{frame}

% \begin{frame}{Research Objectives}
% \begin{enumerate}
% \item \textbf{Compare predictive accuracy} across multiple model architectures
% \item \textbf{Evaluate generalizability} with varying training data sizes
% \item \textbf{Assess robustness} across different time horizons and frequencies
% \item \textbf{Benchmark performance} against naive and ARIMA baselines
% \end{enumerate}
% \end{frame}

% ============================================================
\section{Models}
% ============================================================

\begin{frame}{Models Under Evaluation}
\begin{table}[]
\small
\begin{tabular}{@{}lll@{}}
\toprule
\textbf{Model} & \textbf{Type} & \textbf{Description} \\ \midrule
Naive & Baseline & Simple persistence model using last observation \\
ARIMA & Statistical & Classical auto-regressive model with differencing \\
Chronos & Foundation & Transformer-based pretrained model (Amazon) \\
Sundial & Foundation & Zero-shot forecasting with pretraining \\
TimesFM & Foundation & Google's large-scale time-series foundation model \\ \bottomrule
\end{tabular}
\end{table}
\end{frame}

% ============================================================
\section{Data}
% ============================================================

\begin{frame}{Dataset}
\textbf{Stock:} Microsoft (MSFT)

\vspace{1em}

\textbf{Time Frequencies:}
\begin{itemize}
\item \textbf{Hourly (1h):} High-frequency intraday trading data
\item \textbf{Daily (1d):} End-of-day closing prices
\end{itemize}

\vspace{1em}

\textbf{Test Period:} 2023

\vspace{1em}

\textbf{Training Periods:}
\begin{itemize}
\item Short-term: 25-250 days (linear and logarithmic spacing)
\item Long-term: 1-10 years
\end{itemize}
\end{frame}

\begin{frame}{Data Transformations}
\textbf{Two transformations tested:}
\begin{itemize}
\item \textbf{Raw prices:} Original closing prices
\item \textbf{Percentage change:} Relative change between consecutive observations
\end{itemize}

\vspace{1em}

\textbf{Percentage Change Formula:}
\[
\text{Relative Change} = \frac{\text{New Value} - \text{Old Value}}{\text{Old Value}}
\]

\vspace{1em}

\textbf{Why Percentage Change?}
\begin{itemize}
\item \textbf{Stationarity:} Percentage changes are more stationary than raw prices
\item \textbf{Scale normalization:} Makes different stocks/time periods comparable
\item \textbf{Financial relevance:} Returns are more meaningful than absolute prices for trading decisions
\end{itemize}
\end{frame}

% ============================================================
\section{Experimental Design}
% ============================================================

\begin{frame}{Experimental Design}
\begin{columns}[T]
\column{0.5\textwidth}
\textbf{Training Configurations:}
\begin{itemize}
\item \textbf{Short-term:} 25-250 days
\begin{itemize}
\small
\item Linear spacing
\item Logarithmic spacing
\end{itemize}
\item \textbf{Long-term:} 1-10 years
\end{itemize}

\vspace{1em}

\textbf{Forecast Horizons:}
\begin{itemize}
\item Hourly: 1 hour
\item Daily: 1, 5, 21, 63 days
\end{itemize}

\column{0.5\textwidth}
\textbf{Evaluation Metrics:}
\begin{itemize}
\item MAE (Mean Absolute Error)
\item RMSE (Root Mean Squared Error)
\item MAPE (Mean Absolute \% Error)
\item MSE (Mean Squared Error)
\item MDA (Mean Directional Accuracy)
\end{itemize}

\vspace{1em}

\textbf{Transformations:}
\begin{itemize}
\item Raw prices
\item Percentage change
\end{itemize}
\end{columns}

\vspace{1em}
\begin{block}{Forecasting Approach}
Rolling forecast updating context with actual observations
\end{block}
\end{frame}

% ============================================================
\section{Current Progress}
% ============================================================

\begin{frame}{Experiments Completed}
\begin{itemize}
\item \textbf{All models implemented and tested:}
\begin{itemize}
\item Traditional: ARIMA, Naive baseline
\item Foundation models: Chronos, Sundial, TimesFM
\end{itemize}

\vspace{0.5em}

\item \textbf{Comprehensive experimental setup:}
\begin{itemize}
\item Multiple training configurations (25-250 days, 1-10 years)
\item Multiple forecast horizons (1h, 1d, 5d, 21d, 63d)
\item Two temporal frequencies (hourly, daily)
\item Two data transformations (raw prices, percentage change)
\end{itemize}

\vspace{0.5em}

% \item \textbf{Data collection and processing complete:}
% \begin{itemize}
% \item Full pipeline from data download to evaluation
% \item Automated experiment tracking and configuration management
% \item \textbf{Total:} 5 models × 3 training types × 2 frequencies × 2 transformations = extensive experimental coverage
% \end{itemize}
\end{itemize}

% \vspace{1em}

% \begin{block}{Status}
% Results analysis and interpretation in progress
% \end{block}
\end{frame}

% ============================================================
\section{Next Steps}
% ============================================================

\begin{frame}{Next Steps}
\textbf{RQ2: Temporal Frequency Effects}
\begin{itemize}
\item Extend experiments to minute-level data (higher frequency)
\item Compare model performance across hourly, daily, and minute frequencies
\item Analyze how different temporal resolutions affect foundation vs traditional models
\end{itemize}

\vspace{1em}

\textbf{RQ3: Market Disruption Resilience}
\begin{itemize}
\item Identify high-volatility market regimes (e.g., COVID-19 crash, earnings announcements)
\item Evaluate model performance during market disruptions
\item Compare stability of foundation models vs traditional approaches under stress
\end{itemize}

% \vspace{1em}

% \textbf{Analysis and Writing}
% \begin{itemize}
% \item Complete statistical significance testing
% \item Finalize results interpretation
% \item Thesis writing
% \end{itemize}
\end{frame}

% ============================================================
% \section{Challenges and Limitations}
% ============================================================

% \begin{frame}{Challenges and Limitations}
% \begin{itemize}
% \item \textbf{Data:} Single stock (MSFT), limited market conditions
% \item \textbf{Model constraints:} Conflicting transformer versions (Chronos vs Sundial)
% \item \textbf{Computational:} Foundation models require significant resources
% \item \textbf{Financial markets:} High noise, non-stationarity, external factors
% \end{itemize}

% \vspace{1em}

% \begin{block}{Mitigation Strategies}
% Rigorous experimental design, multiple evaluation metrics, baseline comparisons
% \end{block}
% \end{frame}

% ============================================================
% \section{Timeline}
% ============================================================

% \begin{frame}{Project Timeline}
% \textbf{Completed:}
% \begin{itemize}
% \item Literature review and model selection
% \item Full experimental pipeline implementation
% \item Data collection and preprocessing
% \item All experimental runs (5 models × multiple configurations)
% \end{itemize}

% \vspace{1em}

% \textbf{In Progress:}
% \begin{itemize}
% \item Results analysis and interpretation
% \item Statistical significance testing
% \item Visualization refinement
% \end{itemize}
% \end{frame}

\begin{frame}{Next Steps}
\textbf{Next:}
\begin{itemize}
\item Complete results analysis
\item RQ2: Minute-level data experiments
\item RQ3: High-volatility regime analysis
% \item Thesis writing
\end{itemize}
\end{frame}

% ============================================================
% Conclusion
% ============================================================

% \begin{frame}[standout]
% Thank you!
% \\[1em]
% Questions?
% \end{frame}

\end{document}
